\documentclass[10pt,a4paper]{amsart}
\usepackage[margin=19mm]{geometry}
\usepackage{amsmath,amssymb,mathtools}
\usepackage[scr=rsfs,cal=euler]{mathalfa}
\usepackage{xfrac}
\usepackage[unicode,hidelinks,bookmarks=false]{hyperref}
\newcommand{\ie}{i.e.\ }
\newcommand{\cf}{cf.\ }
\newcommand{\resp}{respectively\ }
\newcommand{\Subsection}{\subsection}
\providecommand{\nobreakdash}{\nobreak\hbox{-}}
\setlength{\emergencystretch}{2em}
\multlinegap0pt
\usepackage{tikz}
\usepackage{enumerate}
\usepackage{amssymb}
\usepackage{tcolorbox}
\newtheorem{prop}{Proposition}
\newtheorem{theorem}{Theorem}
\title{\textbf{A Survey of Baker Wandering Domains}}
\author{Sukanta Das, Tarakanta Nayak}
\date{Source: arXiv:2604.08264v1 (CC BY 4.0)}
\begin{document}
\maketitle

\noindent\emph{Source and licence.} \textbf{A Survey of Baker Wandering Domains}. Version arXiv:2604.08264v1 (CC BY 4.0), distributed by arXiv under the Creative Commons Attribution 4.0 International licence, http://creativecommons.org/licenses/by/4.0/, as recorded in the arXiv licence metadata for this exact version and shown on the abstract page https://arxiv.org/abs/2604.08264v1. Full licence notice: \texttt{licenses/arxiv-2604.08264-CC-BY-4.0.txt}.\par\smallskip

\noindent\textit{Editorial scope.} Counted unit: the article's theorem BWD-LOCALBOV (source0.tex lines 748--751). The article's complete original proof of it is delivered with it (source0.tex lines 752--760, one \texttt{proof} environment of 9 lines). No mathematics is written, repaired or re-derived: the statement and the proof are the article's own lines, in the article's own notation, and no retained line is substituted or re-flowed.  Retained uncounted as the source-local context the proof needs: lines 735--738.  Nothing was omitted from any retained span, so the package carries no omission record.  No retained line contains a citation, so the wrapper carries no bibliography and no bibliography entry.  No retained line contains a figure environment, an image-inclusion command or drawing code, so no image is delivered and the package declares no asset dependency.  Editorial formatting only, in addition: an AMS \texttt{amsart} preamble that restates the article's own declarations and macros used by the retained lines, namely the environments \texttt{prop}, \texttt{theorem} on separate counters, together with the standard packages those lines need.  The retained environments print in this document as Proposition 1, Theorem 1 in order of appearance, whereas the article prints the same items with the numbering of its own full document.  The complete native source of the article as distributed in the arXiv e-print for this exact version is bundled unchanged as \texttt{sources/198146/source0.tex}.
	\begin{prop}
		If $f$ is meromorphic with BWD then the image of every unbounded curve under $f$ is unbounded.  In particular, there is no BWD if the function  has a finite asymptotic value.
		\label{unbounded-BWD}
	\end{prop}

	\begin{theorem}
		If $f$ is meromorphic with at most finitely many poles and has a BWD then it has a local bov at $\infty$.
		\label{BWD-LOCALBOV}
	\end{theorem}

	\begin{proof} Let $D$ be a disk around $\infty$. Then $f^{-1}(D)$ is clearly unbounded (by Picard's theorem). 
		\par
		
		Since each bounded component of $f^{-1}(D)$ has to contain a pole and the number of poles is finite, there is at least one unbounded component of $f^{-1}(D)$. 
		
		\par	
		If the number of unbounded components of $f^{-1}(D)$ is at least two  then one of them has an unbounded boundary component. Let this boundary component be denoted by $\gamma$. This $\gamma$ is mapped into the boundary of $D$,  by $f$. The boundary of $D$ is a bounded set. On the other hand, its image $f(\gamma)$ is unbounded by Proposition~\ref{unbounded-BWD}. This is a contradiction.
		Therefore, the set $f^{-1}(D)$ has a unique  unbounded component, say $C$. Now, using Proposition~\ref{unbounded-BWD} again, we conclude that every component of the boundary of $C$ is  bounded.
	\end{proof} 

\end{document}
